% ===========================================================================
% Section 1 — Introduction
% Contribution claims C1-C5 keep their registered wording (sizing
% included).
% ===========================================================================
\section{Introduction}
\label{sec:introduction}

Instruction in pool and billiards is transmitted through a canon of
empirical rules: the $30^\circ$ and $90^\circ$ rules for post-impact
cue-ball direction, the throw, squirt, and swerve laws, gearing english,
the ten-times-fuller rule for frozen-ball kisses, the trisect system for
draw shots, and the diamond systems for banks and kicks. The canon's
analytic ancestry runs to Coriolis
\citep{coriolis1835,nadler2005coriolis}, and its modern statements are
maintained in Alciatore's technical-proof corpus and the companion
high-speed-video and FAQ collections
\citep{alciatore2008coriolis,alciatore_tp_catalog,alciatore_faq}. Every
rule in the canon is an idealization---a closed form or a geometric
shortcut calibrated at one operating point---yet none is taught with a
quantified statement of \emph{where it holds}. A player learns that a
well-struck draw shot departs at roughly three times the cut angle
``below about $40^\circ$''; no source states the lower boundary of that
claim, its speed window, or the error incurred at the edges.

The literature does not close this gap. Experimental work has validated
single rules at single operating points: \citet{bayes1963collision}
measured ball--ball collision losses, \citet{wallace1988billiard}
analyzed two-dimensional collisions (with the frictional correction of
\citealp{onoda1989comment}), and \citet{cohen2002threecushion} treated
the corner-5 diamond system as geometry alone. Simulation authors have
validated their own sub-models against their own footage
\citep{mathavan2010cushion,mathavan2014numerical}, and the most complete
open simulator, pooltool, ships a single self-consistency notebook
\citep{kiefl2024pooltool}. To our knowledge, no systematic, whole-canon,
domains-of-validity study exists.

This paper conducts that study. The instrument is an event-driven
analytic rigid-body simulator for cue-sports mechanics
(Section~\ref{sec:simulator}): ball motion between events is closed
form, event times are polynomial roots, and the shipped WebAssembly
binary---canonical kernel \code{0.1.0+fb3425ad}---reproduces every
trajectory bit-exactly across the operating systems, architectures, and
runtimes tested (Section~\ref{sec:reproducibility}), making the shipped
binary itself the paper's physics definition
of record. We use the term simulator deliberately: billiards is the
vehicle, and the verification-and-validation (V\&V) methodology together
with the rules study are the contributions. One definition governs the
title's word \emph{validated} and every use of it below: validation means
comparison against anchors classified as \emph{measured data} (camera
points, instrumented tables, high-speed video, a robot protocol), with
comparison error and validation uncertainty reported and every miss
ledgered --- never pass/fail certification; comparisons against sources'
own idealized forms are labeled model benchmarking or verification of the
port, and rule-of-thumb lore is characterized, not validated
(Section~\ref{sec:methodology}). The simulator's validity is
established by a registered battery of 24 experiments spanning
ball--ball impact and throw, stick--ball interaction, cushion rebound
and the diamond systems, curved (Coriolis) trajectories, frozen contacts
and the break, and closed-form kinematics with jump shots, plus a
cross-simulator oracle run---25 pinned outputs in all, byte-stable
across repeated runs.

The battery is designed to be hard on its own instrument
(Section~\ref{sec:methodology}). We adopt ASME V\&V~20 vocabulary
throughout \citep{asme2009vv20}: each anchor names its validation
variable and validation points, reports the comparison error
$E = S - D$ between simulation value $S$ and source datum $D$, and
carries a validation uncertainty $u_{\mathrm{val}}$ whose numerical,
input, and source-data components are estimated per source family
(Section~\ref{sec:methodology}), including digitization uncertainty for
points lifted from published figures.
Verification---agreement of the code with its own closed forms and with
a cross-simulator oracle whose event times match to within
$7.4\times10^{-14}$~s where the model registries coincide---is kept
strictly separate from validation, agreement with physical measurement.
Evidence spans four source tiers (published curves and tables ---
classified per anchor as measured data or derived reference forms ---
high-speed video, a published specification window, and the
cross-simulator oracle), every
experiment design was frozen and registered in-repository before any
sweep ran, and
misses of registered bands are never re-tuned: they are recorded in a
machine-readable waiver ledger (38 rows) kept with the study's artifact
bundle.

\figstub{F0}{F0.pdf}{%
  \textbf{Verdict map.} One row per rule family of
  Table~\ref{tab:domains}, spanning all three anchor classes. Bars
  span the registered sweep domains, verdict-segmented (legend at
  top; in-bar ticks: measured boundaries), with one headline note per
  row (values: pinned outputs, Supplementary Material S1).}{F0}

Figure~\ref{fig:F0} is the battery's one-look summary --- one
verdict-segmented domain bar per rule family, annotated with the
headline comparison errors the following sections quantify. Most rows
benchmark derived forms or lore; ``validation'' is reserved for
measured-datum rows (Section~\ref{sec:methodology}). Throughout the
figures, headline errors $E$ and markers in orange are the
waiver-marker convention (Section~\ref{sec:results}). The map's
239/240 cushion count is verification-of-port; its measured-datum
content is the 12/12 camera points (Section~\ref{sec:results-c}); its
miscue-limit row states the registered input-envelope convention,
revised post-registration under the scoped compatibility note of
Section~\ref{sec:limitations} (simulated results unchanged).
Three numbers preview what the battery finds. The trisect
``$3\times$~cut'' rule is exact only at its half-ball calibration point:
the measured ratio declines monotonically from 3.99 to 2.59 over
$5$--$45^\circ$ cuts and stays within 10\% of 3 only for
$22.5$--$40^\circ$ cuts---even friction-free perfect balls, calibrated
at the same half-ball anchor, give a small-cut limit of 4.0, and the
measured small-cut ratio ($5^\circ$ cut) stays at 3.99--4.01 across the
restitution and friction
sensitivity arms because each arm re-anchors its own calibration; the
breakdown is geometric in that calibration-anchored sense, with $e_b$
and friction setting only the operating point (the decomposition is in
Section~\ref{sec:results-f}) (experiment \expid{F2}; Fig.~\ref{fig:F13}). The
$30^\circ$-rule peak height sits inside the registered
$33.7\pm0.3^\circ$ band only at the slowest registered speed
($E = -0.57^\circ$ at 1.5~m/s and $-0.83^\circ$ at 2.0~m/s, against a
frictionless-elastic reference form), while the peak location and the
$27$--$34^\circ$ plateau hold at every tested speed (experiment
\expid{A2};
Fig.~\ref{fig:F3}). And the standard claim that cut-induced throw
collapses across speeds for cuts at or below $20^\circ$ holds strictly
only to $17.5^\circ$ on the registered grid---consistent with the
source's own crossover near $18^\circ$ (experiment \expid{A3};
Fig.~\ref{fig:F2}).

The methodology also proved to have teeth against the instrument itself:
the battery uncovered two inherited defects in our own kernel---a
ball--ball slip-impulse law that left slip-regime throw about 2.6\% low
(found by experiment \expid{A1}; hereafter the \emph{slip-law defect}) and a
discontinuous settled-rebound branch selection in one
compliant cushion model that jumped registered bank arrivals by
0.49--0.85 diamonds (found by experiment \expid{C1}; hereafter the
\emph{branch-selection defect})---both fixed and
confirmed by predicted-shift verification, and
Section~\ref{sec:oracle} presents the episode as a case study.

\subsection{Contribution claims}
\label{sec:contributions}

This paper makes five contributions, each sized by verified novelty
against the closest prior art we could locate. The novelty sizing rests
on two structured prior-art research sweeps (run 2026-07-21, coverage
through July 2026) across scholar-indexed journals and preprint servers
in the simulation, sports-engineering, and computational-billiards
literature together with the cue-sports gray literature (the Dr.~Dave
technical-proof and high-speed-video catalogs, \emph{Billiards Digest},
the pooltool documentation), with every load-bearing claim verified
against its primary source; the sweep summaries are retained internally
but per-query transcripts were not logged, so we report the search as
scoped rather than as an exhaustively documented systematic-review
protocol, and we claim ``first'' only as first-we-could-locate. The
closest prior art found is listed with each claim.

% Claims C1-C5 keep their registered substance; their "Sizing:"
% annotations are folded into plain prose, and the C4 neighbor check is
% stated as the completed survey sentence.

\begin{description}
  \item[C1.] Systematic domains-of-validity quantification of the empirical
    laws of pool (30$^\circ$/90$^\circ$ rules, throw/squirt/swerve laws,
    gearing english, 10-times-fuller, bank/kick diamond systems, Coriolis
    results) with registered, hash-pinned experiment scripts --- the
    paper's headline
    contribution, addressing a confirmed literature gap. Closest prior
    art: pooltool's single self-consistency
    notebook \citep{kiefl2024pooltool}; \citet{cohen2002threecushion}
    (geometric corner-5 only); \citet{bayes1963collision};
    \citet{wallace1988billiard} (single-rule experiments).
  \item[C2.] Whole-engine multi-source validation battery (published
    curves and tables + HSV footage + published specification windows +
    cross-simulator oracle)
    under an explicit anti-overfit out-of-sample doctrine, with a
    machine-readable waiver ledger --- the paper's headline methodology
    contribution.
    Closest prior art: \citet{mathavan2010cushion,mathavan2014numerical}
    (single sub-models, own footage); no engine-wide precedent.
  \item[C3.] Compliant frozen-cluster break resolution (Hertz +
    Lankarani--Nikravesh, Jacobi concurrent advance) embedded in an
    otherwise closed-form event-driven kernel --- a major section of this
    paper and the seed of a planned companion paper. Closest prior art:
    \citet{gismerosmoreno2022multiple} (full-game penalty time-stepping);
    \citet{cross2023oblique} (small-$N$ analytics); pooltool's 10\%
    momentum-theft hack \citep{kiefl2024pooltool}.
  \item[C4.] The shipped canonical \code{.wasm} as the paper's normative
    physics definition: bit-exact cross-platform reproduction of every
    figure from the shipped binary --- a major section of this paper and,
    to our knowledge, without peer-reviewed simulator precedent. Closest
    prior art: Rapier's enhanced-determinism mode (engineering
    documentation); deterministic-WebAssembly runtimes
    \citep{wasm2015nondeterminism}; the nearest scientific neighbors ---
    versioned WebAssembly tool distributions in bioinformatics and the
    climate community's bit-for-bit reproducibility and
    ensemble-consistency practice --- pin binaries for reuse or test
    statistical consistency of rebuilt source, but neither makes a
    shipped versioned binary a paper's normative physics definition.
  \item[C5.] Terminal rest projection (``round-trip law'': simulate output
    is valid simulate input) and quartic event-detection robustness ---
    deliberately sized as an engineering subsection, not a headline.
    Closest prior art: \citet{leckie2006event} flagged the
    fragility; Algorithm 1010 \citep{orellana2020algorithm1010} + two ACM
    Remarks \citep{demichele2022remark,demichele2024remark}.
\end{description}

The remainder of the paper is organized as follows.
Section~\ref{sec:related} situates the work in the measurement,
simulation, and computational-billiards literature.
Section~\ref{sec:simulator} describes the instrument, and
Section~\ref{sec:methodology} the registered V\&V design, the anchor
tiers, and the waiver-ledger doctrine. Section~\ref{sec:results} reports
the rules study across the six experiment groups, and
Section~\ref{sec:oracle} the cross-simulator agreement study together
with the defect-discovery case study. Section~\ref{sec:reproducibility}
defines the reproduction protocol built on the canonical binary;
Section~\ref{sec:limitations} states limitations, and
Section~\ref{sec:future} outlines future work.
